\section{Introduction}

Les agents intelligents doivent souvent résoudre des problèmes dont la structure
varie fortement : trouver un chemin dans une grille, planifier un itinéraire
routier, réorganiser un puzzle ou construire une expression arithmétique.
Malgré cette diversité, ces tâches peuvent être formulées comme des problèmes
de recherche dans un espace d'états.

L'objectif de ce travail est de concevoir \AUSSI{}, un moteur générique de
résolution de problèmes. Le système doit permettre de résoudre plusieurs types
de problèmes sans modifier les algorithmes d'exploration. Seules la formulation
du problème et, lorsque nécessaire, l'heuristique doivent varier.

Le moteur implémente les stratégies \BFS{}, \UCS{}, \Greedy{} et \AStar{}.
Il permet aussi de choisir entre les modes \TreeSearch{} et \GraphSearch{}.
Les problèmes pris en charge sont les labyrinthes, la navigation dans un graphe
pondéré, le N-Puzzle et le jeu arithmétique.

Les contributions principales du projet sont les suivantes :

\begin{itemize}
    \item Une interface générique \code{Problem} indépendante des domaines.
    \item Un moteur unique \code{SEARCH(problem, strategy, mode)}.
    \item Des stratégies de recherche interchangeables.
    \item Une gestion configurable des états répétés.
    \item Une interface utilisateur permettant de sélectionner un problème,
    une stratégie et un mode de recherche.
    \item Une collecte systématique des métriques d'exécution.
    \item Une étude comparative des algorithmes sur plusieurs instances.
\end{itemize}

Le reste du rapport présente les exigences du projet, l'architecture proposée,
la modélisation des problèmes, les algorithmes implantés, l'interface, les
expériences réalisées et une discussion des résultats.